%%%%%%%%%%%%%%%%%%%%%%%%%%%%%%%%%%%%%%%%%%%%%%%%%%%%%%%%%%%%%
%% Begin exercise %%
%%%%%%%%%%%%%%%%%%%%%%%%%%%%%%%%%%%%%%%%%%%%%%%%%%%%%%%%%%%%%
\ex{Synchronous machines}


\normalsize{\textbf{Acknowledgement}: Parts of the following exercise are adapted from ``Elektrische Antriebstechnik'' by J. Böcker, Paderborn University, 2020
}\\


%%%%%%%%%%%%%%%%%%%%%%%%%%%%%%%%%%%%%%%%%%%%%%%%%%%%%%%%%%%%%
%% Task 1 %%
%%%%%%%%%%%%%%%%%%%%%%%%%%%%%%%%%%%%%%%%%%%%%%%%%%%%%%%%%%%%%

\task{Cylinderical synchronous machine}
In a thermal power station, a cylinderical synchronous generator is used which has the rated data in Tab.~\ref{tab:para_cylindericalSynchonousMachine} and is connected in star. The generator is driven by a turbine. The converted energy is fed into the 50 Hz national grid via a transformer. Saturation effects and losses in the machine can be neglected.
\begin{table}[htb]
    \caption{Parameters of the synchronous machine.}
    \centering
    \begin{tabular}{lll}\toprule
    Symbol  & Description       & Value \\
    \midrule
    $U_{\mathrm{star,n}}$ & Star voltage            & $\SI{\frac{10}{\sqrt{3}}}{\kilo\volt}$ \\
    $I_{\mathrm{star,n}}$ & Star current            & $\SI{6500}{\ampere}$ \\
    $f_{\mathrm{n}}$      & Nominal frequency       & $\SI{50}{\hertz}$ \\
    $\cos(\varphi_{\mathrm{n}})$    & Power factor  & 0.8 (capacitive) \\
    $n_{\mathrm{n}}$        & Nominal speed         & $\SI{1500}{\per\minute}$ \\
    $I_{\mathrm{f,n}}$    & Field current           & $\SI{150}{\ampere}$ \\ 
    $I_{\mathrm{s,sc0}}$  & Short circuit current   & $\SI{7800}{\ampere}$ \\
    \bottomrule
    \end{tabular}
    \label{tab:para_cylindericalSynchonousMachine}
\end{table}




%%%%%%%%%%%%%%%%%%%%%%%%%%%%%%%%%%%%%%%%%%%%%%%%%%%%%%%%%%%%%
\subtask{Synchronization must be carried out to connect the synchronous generator to the electricity grid. Why is this process necessary? Name the synchronization conditions that must be checked before the generator can be connected to the grid.}
\begin{solutionblock}
    Necessity of synchronization:
    \begin{itemize}
        \item Synchronization of the generator with the grid is absolutely essential for a complication-free connection of the machine to the grid. Otherwise, severe transient effects can lead to catastrophic effects.
        \item A high phase shift during the synchronization can start the machine to oscillate, but this is reduced by the damper winding in the rotor.
        \item In addition, connecting the generator with a large phase offset leads to high equalizing currents and thus a large torque are the result, which leads to the destruction of the generator.
    \end{itemize}

    Therefore, the following synchronization conditions should be fulfilled:
    \begin{itemize}
        \item The rotational frequency is equal to the grid frequency.
        \item The terminal voltage and the grid voltage are at the same level.
        \item The sense of rotation of the machine is equal to the grid.
        \item The phase orientation is equal to the grid.
    \end{itemize}

\end{solutionblock}




%%%%%%%%%%%%%%%%%%%%%%%%%%%%%%%%%%%%%%%%%%%%%%%%%%%%%%%%%%%%%
\subtask{How can the amount of active power output be influenced during generator operation? On the other hand, how can the inductive reactive power delivered or absorbed be influenced?}

\begin{solutionblock}
    The delivered active power can only be influenced by changing the mechanical input power. The delivered (over excitation) or absorbed (under excitation) reactive power is controlled via the field current.
\end{solutionblock}



%%%%%%%%%%%%%%%%%%%%%%%%%%%%%%%%%%%%%%%%%%%%%%%%%%%%%%%%%%%%%
\subtask{What is meant by synchronous condenser operation?}

\begin{solutionblock}
    In the synchronous condenser operation mode, only reactive power is delivered into or absorbed from the grid. The machine operates in no-load mode. This operation mode is used to control the reactive power within the grid. Sometimes, more or less reactive power is needed, dependant on the conditions of the considered time frame.
\end{solutionblock}

%%%%%%%%%%%%%%%%%%%%%%%%%%%%%%%%%%%%%%%%%%%%%%%%%%%%%%%%%%%%%
\subtask{Calculate with the given data the pole pair number $p$, the synchronous reactance $X_{\mathrm{s}}$, the apparent power $S$, the phase shift angle $\varphi$, the inner voltage $U_{\mathrm{i}}$ and the load angle $\theta$ at the rated conditions.}

\begin{solutionblock}
    The pole pair number is determined with:
    \begin{equation}
        p = \frac{\omega_{\mathrm{n}}}{\omega_{\mathrm{mech}}}
        = \frac{\SI{2\pi\cdot 50}{\per\second}}{\SI{2\pi \cdot 25}{\per\second}}
        = 2.
    \end{equation}

    The synchronous reactance is calculated by
    \begin{equation}
        X_{\mathrm{s}} = \frac{U_{\mathrm{s}}}{I_{\mathrm{s,sc0}}}
        = \frac{\SI{\frac{10000}{\sqrt{3}}}{\volt}}{\SI{7800}{\ampere}}
        = \SI{0.74}{\Omega},
    \end{equation}
    with the stator voltage $U_{\mathrm{s}}$ and the short circuit current $I_{\mathrm{s,sc0}}$.
    The apparent power determines as
    \begin{equation}
        S = \sqrt{3} U_{\mathrm{n}} I_{\mathrm{n}}
        = \sqrt{3}\cdot \SI{10000}{\volt}\cdot\SI{6500}{\ampere}
        = \SI{112.58}{\mega\volt\ampere},
    \end{equation}
    and this leads to the calculation of the phase shift angle by:
    \begin{equation}
        \varphi = \arccos(-0.8)\cdot \mathrm{sign}\left(\frac{Q}{S}\right)
        = - \arccos(-0.8) = \SI{-143.13}{\degree}.
    \end{equation}

    The inner voltage is calculated with
    \begin{equation}
        U_{\mathrm{i}} = X_{\mathrm{s}} I_{\mathrm{s,sc}}
        = X_{\mathrm{s}} \left(\frac{I_{\mathrm{s,sc0}}-I_{\mathrm{s}}\sin(\varphi)}{\cos(\varphi)}\right),
    \end{equation}
    and the active power is defined as:
    \begin{equation}
        P = \frac{3 U_{\mathrm{s}}U_{\mathrm{i}}}{X_{\mathrm{s}}} \sin(\theta).
    \end{equation}
    Inserting the inner voltage from above into the active power equation, yields:
    \begin{equation}
        P = \frac{3 U_{\mathrm{s}} X_{\mathrm{s}} \left(\frac{I_{\mathrm{s,sc0}}-I_{\mathrm{s}}\sin(\varphi)}{\cos(\varphi)}\right)}{X_{\mathrm{s}}} \sin(\theta)
        = 3 U_{\mathrm{s}} \left(I_{\mathrm{s,sc0}}-I_{\mathrm{s}} \sin(\varphi)\right) \tan(\theta).
    \end{equation}

    Hence, the load angle is determined with resorting the equation from above as follows:
    \begin{align}
        \begin{split}
        \theta &= \arctan\left(\frac{P}{3 U_{\mathrm{s}}(I_{\mathrm{s,sc0}}-I_{\mathrm{s,sc}}\sin(\varphi))}\right)
        =\arctan\left(\frac{3 U_{\mathrm{s}} I_{\mathrm{s}} \cos(\varphi)}{3 U_{\mathrm{s}}(I_{\mathrm{s,sc0}}-I_{\mathrm{s,sc}}\sin(\varphi))}\right) \\
        &= \arctan\left(\frac{I_{\mathrm{s}}\cos(\varphi)}{I_{\mathrm{s,sc0}}-I_{\mathrm{s}}\sin(\varphi)} \right)
        = \arctan\left(\frac{\SI{6.5}{\kilo\ampere}\cdot \cos(\SI{-143.13}{\degree})}{\SI{7.8}{\kilo\ampere}-\SI{6.5}{\kilo\ampere}\cdot\sin(\SI{-143.13}{\degree})}\right)
        = \SI{-23.95}{\degree}.
        \end{split}
        \label{eq:load_angle}
    \end{align}

    The inner voltage is calculated by:
    \begin{equation}
        U_{\mathrm{i}} = X_{\mathrm{s}} I_{\mathrm{s,sc}}
        = X_{\mathrm{s}} \left(\frac{I_{\mathrm{s,sc0}}- I_{\mathrm{s}}\sin(\varphi)}{\cos(\theta)}\right)
        = \SI{0.74}{\Omega}\cdot\left(\frac{\SI{7.8}{\kilo\ampere}-\SI{6.5}{\kilo\ampere}\cdot\sin(\SI{-143.13}{\degree})}{\cos(\SI{-23.95}{\degree})}\right)
        = \SI{9473}{\volt}.
    \end{equation}
\end{solutionblock}


%%%%%%%%%%%%%%%%%%%%%%%%%%%%%%%%%%%%%%%%%%%%%%%%%%%%%%%%%%%%%
\subtask{The generator should only take reactive power from the grid. How large is the stator current $I_{\mathrm{s}}$ and the field current $I_{\mathrm{f}}$ at this operating point, when a reactive power of 120 MVA is extracted from the grid? How large is the phase shift angle $\varphi$ and the load angle $\theta$? Draw the phasor diagram for this operating point.}

\begin{solutionblock}
    The stator current changes to
    \begin{equation}
        I_{\mathrm{s}} = \frac{Q}{\sqrt{3}U_{\mathrm{n}}\sin(\varphi)}
        = \frac{\SI{120}{\mega\volt\ampere}}{\SI{\sqrt{3}\cdot10}{\kilo\volt}}
        = \SI{6928.2}{\ampere},
    \end{equation}

    and the inner voltage is determined with:
    \begin{equation}
        U_{\mathrm{i}} = \frac{3 U_{\mathrm{s}}^2 - X_{\mathrm{s}}Q}{3U_{\mathrm{s}}\cos(\theta)}
        = \frac{3\cdot\left(\frac{\SI{10}{\kilo\volt}}{\sqrt{3}}\right)^2 - \SI{0.74}{\Omega}\cdot\SI{120}{\mega\volt\ampere}}{3\cdot\frac{\SI{10}{\kilo\volt}}{\sqrt{3}}\cos(\SI{0}{\degree})}
        = \SI{1501.11}{\volt}.
    \end{equation}

    Therefore, the necessary field current results in
    \begin{equation}
        I_{\mathrm{f}} = \frac{U_{\mathrm{i}}}{U_{\mathrm{i,n}}}I_{\mathrm{f,n}}
        = \frac{\SI{1501.11}{\volt}}{\SI{9473}{\volt}}\cdot \SI{150}{\ampere}
        = \SI{23.77}{\ampere}.
    \end{equation}

    The phase shift angle is $\varphi = \SI{90}{\degree}$ and the load angle is $\theta = \SI{0}{\degree}$. The phasor diagram is visualized in \autoref{fig:Phasor_diagram_reactive}, the inner voltage $\underline{U}_{\mathrm{i}}$ is purely imaginary, as defined in the lecture notes. The stator voltage $\underline{U}_{\mathrm{s}}$ and the internal voltage are directly on top of each other due to the zero load angle $\theta$.
    \begin{solutionfigure}[ht]
        \centering
        \includegraphics{ex06/Phasor_diagram_reactive.pdf}
        \caption{Phasor diagram for the given operating point. The scala is for the voltage 1 cm $\widehat{=}~ \SI{1000}{\volt}$ and for the current 1 cm $\widehat{=}~ \SI{2000}{\ampere}$.}
        \label{fig:Phasor_diagram_reactive}
    \end{solutionfigure}

\end{solutionblock}


%%%%%%%%%%%%%%%%%%%%%%%%%%%%%%%%%%%%%%%%%%%%%%%%%%%%%%%%%%%%%
\subtask{In contrast to the previous subtask, the generator should provide 80 MW in addition to the reactive power now. How large is the stator current $I_{\mathrm{s}}$, the phase shift angle $\varphi$ and the load angle $\theta$? Draw also for this operation point the phasor diagram.}

\begin{solutionblock}
    
    The apparent power is:
    \begin{equation}
        S = \sqrt{P^2+Q^2}
        = \sqrt{\left(\SI{80}{\mega\watt}\right)^2 + \left(\SI{120}{\mega\volt\ampere}\right)^2}
        = \SI{144.22}{\mega\volt\ampere}.
    \end{equation}

    Hence, the stator current yields:
    \begin{equation}
        I_{\mathrm{s}} = \frac{S}{\sqrt{3}U_{\mathrm{n}}}
        = \frac{\SI{144.22}{\mega\volt\ampere}}{\sqrt{3}\cdot\SI{10}{\kilo\volt}}
        = \SI{8326.54}{\ampere}.
    \end{equation}

    Retrieving the power factor via the complex power equation leads to:
    \begin{equation}
        \varphi = \arccos\left(\frac{-P}{\sqrt{3}U_{\mathrm{s,n}}I_{\mathrm{s}}} \right)
        = \arccos\left(\frac{-\SI{80}{\mega\watt}}{\sqrt{3}\cdot\SI{10}{\kilo\volt}\cdot\SI{8326.54}{\ampere}}\right)
        = \SI{123.7}{\degree}.
    \end{equation}

    The derivation of the load angle is already given in \eqref{eq:load_angle}, which results for the new operation point in
    \begin{equation}
        \theta = \arctan\left(\frac{I_{\mathrm{s}}\cos(\varphi)}{I_{\mathrm{s,sc0}}-I_{\mathrm{s}}\sin(\varphi)}\right)
        = \SI{-79.66}{\degree},
    \end{equation}
    while the resulting inner voltage is
    \begin{equation}
        U_{\mathrm{i}} = X_{\mathrm{s}} I_{\mathrm{s,sc}}
        = \SI{3600}{\volt}.
    \end{equation}

    The phasor diagram is visualized in \autoref{fig:Phasor_diagram_reactive_active}, the inner voltage $\underline{U}_{\mathrm{i}}$ is purely imaginary, as defined in the lecture notes.
    \begin{solutionfigure}
        \centering
        \includegraphics{ex06/Phasor_diagram_reactive_active.pdf}
        \caption{Phasor diagram for the given operating point. The scala is for the voltage 1 cm $\widehat{=}~ \SI{1000}{\volt}$ and for the current 1 cm $\widehat{=}~ \SI{2000}{\ampere}$.}
        \label{fig:Phasor_diagram_reactive_active}
    \end{solutionfigure}


\end{solutionblock}



%%%%%%%%%%%%%%%%%%%%%%%%%%%%%%%%%%%%%%%%%%%%%%%%%%%%%%%%%%%%%
\subtask{Now, the generator should provide a reactive power of 60 MVA into the grid. How large is the generated torque $T$, that the generator operates within the rated apparent power? To which value must the field current $I_{\mathrm{f}}$ change, such that the rated current is reached in the stator winding?}

\begin{solutionblock}
    With the determined rated apparent power in the task before, the given reactive power, the maximum power is calculated by
    \begin{equation}
        P = \sqrt{S_{\mathrm{n}}^2 - Q^2}
        = \sqrt{\left(\SI{112.58}{\mega\volt\ampere}\right)^2 - \left(\SI{60}{\mega\volt\ampere}\right)^2}
        = \SI{95.26}{\mega\watt},
    \end{equation}
    hence, the maximum torque is given with:
    \begin{equation}
        T_{\mathrm{max}} = \frac{P}{\omega_{\mathrm{mech,n}}}
        = \frac{\SI{95.26}{\mega\watt}}{\SI{2\pi\cdot 25}{\per\second}}
        = \SI{606.5}{\kilo\newton\metre}.
    \end{equation}
    

    The angle $\varphi$ for the phase shift at nominal current is determined as:
    \begin{equation}
        \varphi = \arccos\left(\frac{P}{\sqrt{3}U_{\mathrm{star,n} I_{\mathrm{star,n}}}}\right)
        = \arccos\left(\frac{-\SI{95.26}{\mega\watt}}{\sqrt{3}\cdot \SI{10}{\kilo\volt}\cdot \SI{6.5}{\kilo\ampere}}\right)
        = -\SI{147.8}{\degree}.
    \end{equation}

    With \eqref{eq:load_angle} the load angle is calculated by
    \begin{equation}
        \theta = \arctan\left(\frac{I_{\mathrm{s}}\cos(\varphi)}{I_{\mathrm{s,sc0}}-I_{\mathrm{s,sc}}\sin(\varphi)}\right)
        = \arctan\left(\frac{\SI{6.5}{\kilo\ampere}\cos(\SI{-147.8}{\degree})}{\SI{7.8}{\kilo\ampere}-\SI{6.5}{\kilo\ampere}\sin(\SI{-147.8}{\degree})}\right)
        = \SI{-26.03}{\degree},
    \end{equation}
    therefore, the inner voltage results into:
    \begin{equation}
        U_{\mathrm{i}} = X_{\mathrm{s}} I_{\mathrm{s,sc}}
        = X_{\mathrm{s}} \left(\frac{I_{\mathrm{s,sc0}}-I_{\mathrm{s}}\sin(\varphi)}{\cos(\theta)}\right)
        = \SI{0.74}{\Omega}\left(\frac{\SI{7.8}{\kilo\ampere}-\SI{6.5}{\kilo\ampere}\sin(\SI{-147.8}{\degree})}{\cos(\SI{-26.03}{\degree})}\right)
        = \SI{9276}{\volt}.
    \end{equation}

    The field current must be set to the following value:
    \begin{equation}
        I_{\mathrm{f}} = \frac{U_{\mathrm{i}}}{U_{\mathrm{i,n}}}I_{\mathrm{f,n}}
        = \frac{\SI{9276}{\volt}}{\SI{9473}{\volt}}\cdot \SI{150}{\ampere}
        = \SI{146.9}{\ampere}.
    \end{equation}

\end{solutionblock}

\FloatBarrier


%%%%%%%%%%%%%%%%%%%%%%%%%%%%%%%%%%%%%%%%%%%%%%%%%%%%%%%%%%%%%
%% Task 2 %%
%%%%%%%%%%%%%%%%%%%%%%%%%%%%%%%%%%%%%%%%%%%%%%%%%%%%%%%%%%%%%

\task{Steady-state operation of a synchronous machine}
Consider a loss-free synchronous machine with the parameters from Tab.~\ref{tab:para_SynchonousMachine}. First, the synchronous machine is operated as a a motor. Using the data given above, determine the following characteristic variables of the synchronous machine at rated operation. An example phasor diagram for this operation is shown in \autoref{fig:Phasor_diagrams_m_oe}.


\begin{table}[htb]
    \caption{Parameters of the synchronous machine.}
    \centering
    \begin{tabular}{lll}\toprule
    Symbol  & Description       & Value \\
    \midrule
    $U_{\mathrm{star,n}}$ & Star voltage            & $\SI{\frac{6}{\sqrt{3}}}{\kilo\volt}$ \\
    $I_{\mathrm{star,n}}$ & Star current            & $\SI{96}{\ampere}$ \\
    $f_{\mathrm{n}}$      & Nominal frequency       & $\SI{50}{\hertz}$ \\
    $p$                   & Pole pair number        & 2 \\
    $\cos(\varphi_{\mathrm{n}})$    & Power factor  & 0.9 (capacitive) \\
    $T_{\mathrm{n}}$      & Nominal torque          & $\frac{T_{\mathrm{max}}}{2}$ \\
    \bottomrule
    \end{tabular}
    \label{tab:para_SynchonousMachine}
\end{table}

\begin{figure}[ht]
    \centering
    \includegraphics{ex06/Phasor_diagrams_m_oe.pdf}
    \caption{Phasor diagram of a over excited synchronous machine in motor mode.}
    \label{fig:Phasor_diagrams_m_oe}
\end{figure}

%%%%%%%%%%%%%%%%%%%%%%%%%%%%%%%%%%%%%%%%%%%%%%%%%%%%%%%%%%%%%
\subtask{How large is the maximum torque $T_{\mathrm{max}}$?}

\begin{solutionblock}

    Due to the ideal machine characteristic (no losses), the nominal torque is calculated by
    \begin{equation}
        T_{\mathrm{n}} = \frac{P_{\mathrm{n}}}{\omega_{\mathrm{mech,n}}}
        = \frac{3 U_{\mathrm{star,n}} I_{\mathrm{star,n}} \cos(\varphi)}{\omega_{\mathrm{mech,n}}}
        = \frac{\sqrt{3}\cdot \SI{6000}{\volt}\cdot\SI{96}{\ampere}\cdot0.9}{\SI{2\pi\cdot50}{\per\second}}
        = \SI{2858}{\newton\metre},
    \end{equation}

    and, the maximum torque calculates as follows:
    \begin{equation}
        T_{\mathrm{max}} = 2 T_{\mathrm{n}}
        = 2 \cdot \SI{2858}{\newton\metre}
        = \SI{5716}{\newton\metre}.
    \end{equation}
    
\end{solutionblock}



%%%%%%%%%%%%%%%%%%%%%%%%%%%%%%%%%%%%%%%%%%%%%%%%%%%%%%%%%%%%%
\subtask{Which load angle $\theta$ is set at the nominal operating point?}

\begin{solutionblock}
    The load angle can be determined by geometric transformations, which are shown in \autoref{fig:syn_motor_overExcited}.
    \begin{solutionfigure}
        \centering
        \includegraphics{ex06/syn_motor_overExcited.pdf}
        \caption{Phasor diagram for an over excited synchronous machine in motor operation.}
        \label{fig:syn_motor_overExcited}
    \end{solutionfigure}
    
    The target is to create two different triangles, one with the known angle $\varphi$ (marked in green in the figure) and one without the load angle $\theta$ (marked in gray in the figure). 
    Hence, the load angle calculates as follows
    \begin{align}
        \begin{split}
            \sin(\theta) I_{\mathrm{s,sc}} = \cos(\varphi)I_{\mathrm{s}},\\
            \sin(\theta) = \frac{I_{\mathrm{s}}}{I_{\mathrm{s,sc}}}\cos(\varphi),
        \end{split}
    \end{align}
    which leads to the angle:
    \begin{equation}
        \theta = \arcsin\left(\frac{I_{\mathrm{s}}}{I_{\mathrm{s,sc}}} \right) \cos(\varphi)
        = \arcsin\left(\frac{\SI{96}{\ampere}}{\SI{172.8}{\ampere}}\right) \cdot 0.9
        = \SI{30}{\degree}.
    \end{equation}
\end{solutionblock}

%%%%%%%%%%%%%%%%%%%%%%%%%%%%%%%%%%%%%%%%%%%%%%%%%%%%%%%%%%%%%
\subtask{Determine the synchronous reactance $X_{\mathrm{s}}$ value.}

\begin{solutionblock}
    The reactance is determined as:
    \begin{align}
        \begin{split}
        X_{\mathrm{s}}
        &= \frac{U_{\mathrm{i}}}{I_{\mathrm{s,sc}}}
        = \frac{U_{\mathrm{s}}}{I_{\mathrm{s,sc0}}}
        = \frac{U_{\mathrm{s}}}{I_{\mathrm{s}}\sin(\arccos(0.9)\mathrm{sign}(Q/S))+I_{\mathrm{s,sc}}\cos(\theta)} \\
        &= \frac{\SI{\frac{6000}{\sqrt{3}}}{\volt}}{\SI{96}{\ampere}\cdot \sin(\SI{-25.8}{\degree}) + \SI{172.8}{\ampere}\cdot \cos(\SI{30}{\degree})}
        = \SI{32.12}{\Omega}.
        \end{split}
    \end{align}

\end{solutionblock}


%%%%%%%%%%%%%%%%%%%%%%%%%%%%%%%%%%%%%%%%%%%%%%%%%%%%%%%%%%%%%
\subtask{Calculate the inner voltage $U_{\mathrm{i}}$ at the nominal operating point.}
\begin{solutionblock}
    
    \begin{equation}
        U_{\mathrm{i}} = X_{\mathrm{s}} I_{\mathrm{s,sc}}
        = \SI{32.12}{\Omega} \cdot \SI{172.79}{\ampere}
        = \SI{5549.52}{\volt}
    \end{equation}
\end{solutionblock}




%%%%%%%%%%%%%%%%%%%%%%%%%%%%%%%%%%%%%%%%%%%%%%%%%%%%%%%%%%%%%
\subtask{In the next task, the machine operates as generator and the power is 500 kW. The phasor diagram for this operation is shown in \autoref{fig:Phasor_diagrams_gen_oe}. Which load angle $\theta$ is set at the new operating point?}

\begin{figure}[ht]
    \centering
    \includegraphics{ex06/Phasor_diagrams_gen_oe.pdf}
    \caption{Phasor diagram of an over excited synchronous machine in generator mode.}
    \label{fig:Phasor_diagrams_gen_oe}
\end{figure}


\begin{solutionblock}
    In the generator mode, the load angle is negative. With the generated power of 500 kW and the constant voltage, the product of the current and power factor is determined by:
    \begin{equation}
        \cos(\varphi) I_{\mathrm{s}} = \frac{P}{3 U_{\mathrm{s}}}
        = \frac{\SI{-500}{\kilo\watt}}{3\cdot \SI{\frac{6000}{\sqrt{3}}}{\volt}}
        = -\SI{48.11}{\ampere}.
    \end{equation}

    It is also assumed that the short-circuit current has not changed its value by maintaining the excitation current. The load angle can now be determined using the trigonometric relationship:
    \begin{equation}
        \theta = \arcsin\left(\frac{I_{\mathrm{s}}}{I_{\mathrm{s,sc}}}\cos(\varphi)\right)
        = \arcsin\left(\frac{\SI{-48.11}{\ampere}}{\SI{172.8}{\ampere}}\right)
        = \SI{-16.17}{\degree}.
    \end{equation}
\end{solutionblock}


%%%%%%%%%%%%%%%%%%%%%%%%%%%%%%%%%%%%%%%%%%%%%%%%%%%%%%%%%%%%%
\subtask{Calculate the stator current $I_{\mathrm{s}}$.}

\begin{solutionblock}
    The stator current is calculated with the active and the reactive current components, which are defined as in the previous task by
    \begin{equation}
        I_{\mathrm{s}} = \sqrt{I_{\mathrm{active}}^2 + I_{\mathrm{reactive}}^2}
        = \sqrt{\left(I_{\mathrm{s}}\sin(-\varphi - \SI{90}{\degree})\right)^2 + \left(I_{\mathrm{s}}\cos(-\varphi - \SI{90}{\degree})\right)^2}.
    \end{equation}
    
    With the relationship between the trigonometric functions and the phase shift, the equation from above is resorted to:
    \begin{align}
        \begin{split}
            I_{\mathrm{s}} &= \sqrt{\left(-I_{\mathrm{s}}\cos(-\varphi)\right)^2 + \left(I_{\mathrm{s}}\sin(-\varphi)\right)^2} \\
            & = \sqrt{\left(-I_{\mathrm{s}}\cos(-\varphi)\right)^2 + \left(-I_{\mathrm{s}}\sin(\varphi)\right)^2}.
        \end{split}
    \end{align}

    With the triangular relationship:
    \begin{align}
        \begin{split}
        I_{\mathrm{s}} &= \sqrt{\left(-I_{\mathrm{s}}\cos(-\varphi)\right)^2 + \left(I_{\mathrm{s,sc0}}-I_{\mathrm{s,sc}}\cos(\theta)\right)^2} \\
        &= \sqrt{\left(-I_{\mathrm{s}}\cos(-\varphi)\right)^2 + \left(I_{\mathrm{s,sc0}}+I_{\mathrm{s,sc}}\cos(\theta)\right)^2}.
        \end{split}
    \end{align}

    Use the calculation of the short circuit current $I_{\mathrm{s,sc0}}$, leads to:
    \begin{align}
        \begin{split}
            I_{\mathrm{s}} &= \sqrt{\left(-I_{\mathrm{s}}\cos(-\varphi)\right)^2 + \left(-\frac{U_{\mathrm{s}}}{X_{\mathrm{s}}}+I_{\mathrm{s,sc}}\cos(\theta)\right)^2} \\
            &= \sqrt{\left(\SI{48.11}{\ampere}\right)^2 + \left(-\frac{\SI{\frac{6000}{\sqrt{3}}}{\volt}}{\SI{32.12}{\Omega}} + \SI{172.79}{\ampere}\cos(-\SI{16.17}{\degree})\right)^2} \\
            &= \SI{75.44}{\ampere}.
        \end{split}
    \end{align}

\end{solutionblock}


%%%%%%%%%%%%%%%%%%%%%%%%%%%%%%%%%%%%%%%%%%%%%%%%%%%%%%%%%%%%%
\subtask{How large is the power factor $\cos(\varphi)$ at this operating point? In addition, calculate the angle $\varphi$.}

\begin{solutionblock}
    The active power is calculated with
    \begin{equation}
        I_{\mathrm{active}} = I_{\mathrm{s}} \cos(\varphi),
    \end{equation}
    which than leads to
    \begin{equation}
        \cos(\varphi) = \frac{I_{\mathrm{active}}}{I_{\mathrm{s}}} 
        = \frac{I_{\mathrm{s}}\cos(\varphi)}{I_{\mathrm{s}}}
        = \frac{-\SI{48.11}{\ampere}}{\SI{75.44}{\ampere}}
        = -0.638,
    \end{equation}
    and, therefore, the angle results into:
    \begin{equation}
        \varphi = \arccos(-0.638)\cdot \mathrm{sign}\left(\frac{Q}{S}\right)
        = -\arccos(-0.638) = \SI{-129.6}{\degree}.
    \end{equation}
\end{solutionblock}








